\setchapterpreamble[u]{\margintoc}
\chapter{O SQL Server e o Banco de Dados}
\labch{cet615-sqlserver}

Este capítulo cobre a segunda apresentação da primeira
aula~\cite{BaiocoCET615}: a ferramenta de administração, os bancos que o
próprio SQL Server mantém, os objetos de um banco e os arquivos em disco.
Os complementos vêm da documentação oficial no Microsoft Learn.

\section{SQL Server Management Studio}

O \textbf{SSMS}\index{SSMS} é a interface gráfica para administrar o SQL
Server desde a versão 2005; antes dela, o papel era do \emph{Enterprise
Manager}. Quase tudo o que se faz em T-SQL também pode ser feito pelo SSMS,
e o editor de consultas do SSMS executa os comandos T-SQL diretamente.

\section{Bancos do sistema e bancos de usuário}

Um servidor SQL Server guarda dois tipos de banco: os \textbf{bancos do
sistema}, que o SGBD usa para se controlar, e os \textbf{bancos de
usuário}, criados para as aplicações. Os bancos do sistema guardam
metadados, os ``dados sobre os dados''.

\begin{description}
	\item[\texttt{master}:]\index{master@\texttt{master}} registra o nível de
	sistema da instância: logins, bancos existentes, espaço de cada um,
	configurações. A aula recomenda manter sempre um backup atualizado dele.
	\item[\texttt{model}:] serve de molde para os bancos novos. Tudo o que
	estiver no \texttt{model} (opções, tamanhos, permissões padrão, tipos)
	é copiado para cada banco criado depois.
	\item[\texttt{tempdb}:] área de trabalho para tabelas temporárias e
	resultados intermediários de consultas, como as que usam
	\texttt{group by} e \texttt{order by}.
	\item[\texttt{msdb}:] usado pelo SQL Server Agent para agendar tarefas
	(\emph{jobs}) e alertas.
\end{description}

Em aula\sidecite{BaiocoCET615Aulas}, a professora reforçou que o SQL
Server \emph{distribui} o catálogo: o que vale para a instância toda
(logins, lista de bancos) fica no \texttt{master}, e o esquema de cada
banco fica nas tabelas de sistema do próprio banco. Sem o \texttt{master},
os arquivos no disco não podem ser interpretados. Essa divisão é do SQL
Server; outros SGBDs, como o Oracle, organizam o catálogo de outro jeito.

Ela mostrou isso ao vivo: \texttt{select * from sys.databases}, executado
no \texttt{master}, devolve a mesma lista de bancos que o SSMS exibe, com
mais colunas. Todo comando DDL apenas grava no catálogo; depois de um
\texttt{create database}, aparece uma linha nova em \texttt{sys.databases}.
O SSMS é ``uma casca'' sobre esse catálogo, e o banco continua existindo
sem a ferramenta. Um efeito prático: o SSMS só relê o catálogo quando se
pede, então um banco criado por script só aparece no \emph{Object
Explorer} depois de um \emph{Refresh}.

\subsection*{Complemento: o banco \texttt{Resource} e a vida do \texttt{tempdb}}

A documentação lista um quinto banco do sistema que não aparece no SSMS:
o \texttt{Resource}\sidecite{MSSystemDatabases}. Ele é somente leitura e
guarda fisicamente os objetos de sistema que acompanham o produto; esses
objetos aparecem, logicamente, no esquema \texttt{sys} de cada banco.

Sobre o \texttt{tempdb}, a documentação\sidecite{MSTempdb} detalha dois
momentos de limpeza diferentes:

\begin{itemize}
	\item tabelas temporárias locais (\texttt{\#tabela}) e procedimentos
	temporários somem quando a sessão que os criou se desconecta;
	\item o \texttt{tempdb} inteiro é recriado vazio toda vez que o serviço
	do SQL Server inicia. Por isso não se faz backup dele.
\end{itemize}

Além das tabelas temporárias, o \texttt{tempdb} guarda objetos internos,
como as tabelas de trabalho de ordenações e agrupamentos, e o
\emph{version store}, usado pelo isolamento por versão de linha e pelos
gatilhos \texttt{after} (\refch{cet615-gatilhos}).

\section{Tabelas do sistema e procedimentos do sistema}

Os slides listam tabelas de sistema presentes em cada banco
(\texttt{sysobjects}, \texttt{syscolumns}, \texttt{sysindexes},
\texttt{sysforeignkeys}, \texttt{systypes}, \texttt{sysusers}) e outras só
no \texttt{master} (\texttt{sysdatabases}, \texttt{syslogins},
\texttt{sysmessages}, \texttt{sysprocesses}, \texttt{sysconfigures}).

Os \textbf{procedimentos do sistema}\index{procedimento do sistema}, com
prefixo \texttt{sp\_}, automatizam tarefas comuns:

\begin{lstlisting}
exec sp_help obras_literarias   -- estrutura de um objeto
exec sp_helpdb espbd            -- dados de um banco e seus arquivos
exec sp_helpindex compra        -- índices de uma tabela
exec sp_helptext ins_livro      -- código de visão, procedimento ou gatilho
\end{lstlisting}

\subsection*{Complemento: visões de compatibilidade e visões de catálogo}

Desde o SQL Server 2005, as tabelas \texttt{sys...} do SQL Server 2000
existem só como \textbf{visões de compatibilidade}\sidecite{MSCompatViews}.
Elas mostram apenas os metadados que existiam na versão 2000 e podem
devolver nulo ou estourar em colunas de identificadores. Para recursos
novos, a documentação recomenda as \textbf{visões de catálogo}\index{visão
de catálogo}:

\begin{table}[h]
	\begin{tabular}{ll}
		\toprule
		Compatibilidade & Catálogo \\
		\midrule
		\texttt{sysobjects} & \texttt{sys.objects} \\
		\texttt{syscolumns} & \texttt{sys.columns} \\
		\texttt{sysindexes} & \texttt{sys.indexes} \\
		\texttt{sysforeignkeys} & \texttt{sys.foreign\_keys} \\
		\texttt{systypes} & \texttt{sys.types} \\
		\texttt{sysusers} & \texttt{sys.database\_principals} \\
		\texttt{sysdatabases} & \texttt{sys.databases} \\
		\texttt{sysprocesses} & \texttt{sys.dm\_exec\_requests} \\
		\bottomrule
	\end{tabular}
	\caption{Visões de compatibilidade e as visões de catálogo que as substituem.}
\end{table}

A documentação também desaconselha alterar objetos de sistema
diretamente: a administração deve passar pelo SSMS, pela DDL e pelos
procedimentos do sistema\sidecite{MSSystemDatabases}.

\section{Objetos de um banco}

Cada banco contém objetos (componentes lógicos):

\begin{description}
	\item[Tabelas:] as relações do modelo relacional.
	\item[Índices:] aceleram o acesso às linhas das tabelas.
	\item[Visões:] uma forma alternativa de ver dados de uma ou mais
	tabelas (\refch{cet615-visoes}).
	\item[Procedimentos armazenados:] conjuntos de comandos SQL compilados
	e guardados no servidor (\refch{cet615-procedimentos}).
	\item[Gatilhos:] procedimentos especiais que o SGBD executa sozinho
	quando há \texttt{insert}, \texttt{update} ou \texttt{delete} na tabela
	a que estão ligados (\refch{cet615-gatilhos}).
	\item[Outros:] tipos de dados, restrições, \emph{defaults} e regras.
\end{description}

\section{T-SQL}

O \textbf{Transact-SQL}\index{T-SQL} é a implementação da SQL no SQL
Server: segue o padrão ANSI/ISO e acrescenta extensões, como variáveis,
controle de fluxo e procedimentos. A divisão da aula para o T-SQL é:

\begin{description}
	\item[DDL:] criar, alterar e remover bancos, tabelas, índices, visões,
	procedimentos, gatilhos, tipos, restrições, \emph{defaults} e regras.
	\item[DML:] \texttt{select}, \texttt{insert}, \texttt{update},
	\texttt{delete}, transações e a execução de visões e procedimentos.
	\item[DCL:] logins, usuários e grupos, privilégios, backup e
	recuperação, \emph{jobs}.
\end{description}

Cada fabricante tem sua extensão procedural: T-SQL no SQL Server, PL/SQL
no Oracle, SQL PL no DB2. A SQL é declarativa (diz \emph{o que} se quer);
os comandos procedurais são imperativos (dizem \emph{como} fazer). Mesmo
comandos padronizados variam entre fabricantes, por exemplo nos tipos de
dados. O \texttt{create trigger} existe no padrão, mas o corpo do gatilho
é escrito na extensão de cada um.

\subsection*{Complemento: lotes e o comando \texttt{go}}

\index{go@\texttt{go}}Os scripts da disciplina separam trechos com
\texttt{go}. Em aula, a professora disse que ele não é necessário no SSMS
e que existe para utilitários de linha de comando. A documentação
precisa esse ponto\sidecite{MSGo}: \texttt{go} não é um comando T-SQL e
nunca chega ao servidor, mas é reconhecido tanto pelo \texttt{sqlcmd}
quanto pelo editor do SSMS. Ele marca o fim de um \textbf{lote}
(\emph{batch}): tudo o que vem antes é enviado e compilado junto.

Por isso o \texttt{go} faz diferença mesmo no SSMS quando se executa um
script inteiro:

\begin{itemize}
	\item variáveis locais só existem dentro do lote em que foram
	declaradas;
	\item \texttt{create view}, \texttt{create procedure} e
	\texttt{create trigger} precisam ser o primeiro comando do lote;
	\item fora da primeira linha do lote, a chamada de um procedimento
	precisa de \texttt{exec}.
\end{itemize}

Ao executar só um trecho selecionado, como a professora faz em aula, o
trecho já é um lote, e o \texttt{go} não muda nada.

\section{Arquivos de um banco}

Cada banco corresponde a arquivos no disco:

\begin{description}
	\item[Arquivo de dados primário (\texttt{.mdf}):] contém as informações
	de inicialização do banco, aponta para os demais arquivos e guarda as
	tabelas do sistema. Todo banco tem exatamente um.
	\item[Arquivos de dados secundários (\texttt{.ndf}):] opcionais;
	acrescentados conforme o banco cresce.
	\item[Arquivo de log de transações (\texttt{.ldf}):] registra as
	operações feitas no banco. Pode haver mais de um. A regra prática da
	aula é dimensionar o log entre 10\% e 15\% do tamanho do banco (em
	aula, ela falou em até 20\%).
\end{description}

O log não acompanha o tamanho do banco porque cresce e encolhe conforme
as transações terminam e o backup de log é feito. Um log que cresce
demais indica transações longas ou backup de log que não está sendo
feito. A professora também chamou a atenção para a separação entre a
visão lógica e a física: o usuário vê tabelas, mas no disco há poucos
arquivos, não um por tabela. Pelo arquivo não se descobre a estrutura do
banco nem se ele está cheio, porque o tamanho é reservado antes. Isso
soma à segurança, junto com a criptografia.

\subsection*{Complemento: grupos de arquivos}

A documentação\sidecite{MSFilesFilegroups} acrescenta pontos que a aula
não trata:

\begin{itemize}
	\item as extensões \texttt{.mdf}, \texttt{.ndf} e \texttt{.ldf} são
	recomendadas, não obrigatórias;
	\item arquivos de dados se organizam em \textbf{grupos de arquivos}
	(\emph{filegroups})\index{filegroup}. O grupo \texttt{PRIMARY} contém o
	arquivo primário e é o grupo padrão, a menos que outro seja definido
	com \texttt{alter database}. Uma tabela pode ser criada em um grupo
	específico (\texttt{create table ... on grupo});
	\item arquivos de log nunca pertencem a grupo de arquivos;
	\item dentro de um grupo, o SQL Server distribui as gravações entre os
	arquivos em proporção ao espaço livre de cada um;
	\item a recomendação é separar dados e log em discos diferentes; a
	maioria dos bancos funciona bem com um arquivo de dados e um de log.
\end{itemize}

Cada arquivo tem dois nomes: o \textbf{lógico} (\texttt{name}), usado nos
comandos T-SQL, e o \textbf{físico} (\texttt{filename}), o caminho no
sistema operacional.

\section{Criar, alterar e excluir bancos}

Os três comandos são executados a partir do banco \texttt{master}.

\subsection{\texttt{create database}}

\begin{lstlisting}[caption={Banco com um arquivo de dados e um de log.}]
use master
go

create database biblioteca
on
   (name = biblioteca_dados,
    filename = 'C:\dados\biblioteca.mdf',
    size = 10MB,
    maxsize = 50MB,
    filegrowth = 5MB)
log on
   (name = biblioteca_log,
    filename = 'C:\dados\biblioteca.ldf',
    size = 2MB,
    maxsize = 10MB,
    filegrowth = 10%)
\end{lstlisting}

Observações da aula sobre o comando:

\begin{itemize}
	\item \texttt{size} e \texttt{maxsize} aceitam KB, MB, GB e TB; sem
	unidade, vale MB;
	\item \texttt{filegrowth} pode ser absoluto (\texttt{5MB}) ou
	percentual (\texttt{10\%});
	\item não se usam decimais: em vez de \texttt{2.5MB}, escreve-se
	\texttt{2560KB};
	\item uma instância aceita no máximo 32.767 bancos;
	\item o caminho em \texttt{filename} depende da instalação; o script
	da aula usa a pasta \texttt{DATA} da instância.
\end{itemize}

\subsection*{Complemento: o que é obrigatório no \texttt{create database}}

Em aula, a professora apresentou o \texttt{log on} como obrigatório e o
arquivo secundário como opcional (normalmente acrescentado depois, com
\texttt{alter database}). A documentação\sidecite{MSCreateDatabase} diz
que todo banco tem pelo menos dois arquivos, um de dados primário e um
de log, mas a \emph{cláusula} \texttt{log on} é opcional: sem ela, o SQL
Server cria sozinho um arquivo de log com 25\% da soma dos arquivos de
dados ou 512 KB, o que for maior. Sem nenhuma especificação de arquivo,
o primário recebe o tamanho do primário do \texttt{model}; sem
\texttt{maxsize}, o arquivo cresce até ocupar o disco. O
\texttt{create database} também não pode rodar dentro de uma transação
explícita.

\begin{lstlisting}
create database rascunho   -- tamanhos e caminhos padrão
\end{lstlisting}

\subsection{\texttt{alter database}}

Todo \texttt{alter} pressupõe um banco em uso: é comando de manutenção,
corretiva ou por mudança de regra de negócio. Em desenvolvimento, apagar
e recriar costuma ser mais simples. Remover arquivo ou reduzir tamanho
pode perder dados, então a recomendação da aula é fazer backup antes.

\begin{lstlisting}[caption={Alterações comuns nos arquivos de um banco.}]
-- acrescentar um arquivo de dados secundário
alter database biblioteca
add file
   (name = biblioteca_dados2,
    filename = 'C:\dados\biblioteca2.ndf',
    size = 5MB)

-- aumentar o tamanho de um arquivo
alter database biblioteca
modify file (name = biblioteca_dados, size = 20MB)

-- remover um arquivo (ele precisa estar vazio)
alter database biblioteca
remove file biblioteca_dados2
\end{lstlisting}

\subsection{\texttt{drop database}}

Nas palavras da professora, o \texttt{drop} é o comando mais simples e o
mais perigoso, para qualquer objeto.

\begin{lstlisting}
drop database biblioteca
drop database banco1, banco2   -- vários de uma vez
\end{lstlisting}

\subsection{Pelo SSMS}

\begin{itemize}
	\item criar: botão direito em \emph{Databases} $\rightarrow$ \emph{New
	Database};
	\item alterar: botão direito no banco $\rightarrow$ \emph{Properties};
	\item excluir: botão direito no banco $\rightarrow$ \emph{Delete}.
\end{itemize}
